% GENERATED FILE. Edit canonical lessons, then run npm run generate:books.
% canonical-source-hash: fnv1a64:bb2a254aed2ef4b2
% canonical-lessons: MW-C274-takliph, MW-C274-pukar, MW-C274-saharo, MW-C274-daya, MW-C274-bhikh

\chapter{Help in Trouble}
\label{ch:mw-help-in-trouble}
\hlchaptermodality{274}

\begin{chapteropening}
\textbf{By the end of this chapter:} \emph{I can say that someone is in trouble and needs help.}

Complete the last lesson of chapter 274: I can say that someone is in trouble and needs help.
\end{chapteropening}

\section[\texorpdfstring{\mw{तकलीफ} (taklīph)}{तकलीफ (taklīph)}]{\mw{तकलीफ} (taklīph) — trouble}

\label{lesson:MW-C274-takliph}



\begin{quote}
Before the new one: say the Marwadi for to keep secret, then the Marwadi for a promise.
\end{quote}

\subsection*{You'll want to know: \mw{तकलीफ}}
\textbf{\mw{तकलीफ}} — \emph{taklīph} — \textquotedblleft{}trouble\textquotedblright{}.

\begin{quote}
\mw{कोई} \mw{तकलीफ} \mw{है}?
\end{quote}

\begin{quote}
\emph{koī taklīph hai?}
\end{quote}

\begin{quote}
Is anything the matter?
\end{quote}

A feminine noun.

\begin{cousinweb}
\textbf{\mw{तकलीफ}} (\emph{taklīph}) is usually traced, through Persian \emph{taklīf}, to Arabic \emph{taklīf}, \textquotedblleft{}a burden laid on someone\textquotedblright{}.
\end{cousinweb}

\begin{grammarlens}[title={What you've built}]
One more word for asking and answering.
\end{grammarlens}

\subsection*{Your turn}
\begin{itemize}
\raggedright
  \item \emph{Say it:} \emph{taklīph}
  \item \emph{Say it:} \emph{taklīph}, once more
  \item \emph{Say it:} \emph{koī taklīph hai?} — Is anything the matter?
  \item \emph{Recall:} say the Marwadi for to keep secret, then the Marwadi for a promise, then say \emph{taklīph} again
\end{itemize}

\begin{tcolorbox}[breakable,colback=teal!4,colframe=teal!35!black,title={Before you move on}]
What does \mw{तकलीफ} mean? (\textquotedblleft{}Trouble\textquotedblright{}.) Say it once more.
\end{tcolorbox}
\section[\texorpdfstring{\mw{पुकार} (pukār)}{पुकार (pukār)}]{\mw{पुकार} (pukār) — a cry for help}

\label{lesson:MW-C274-pukar}



\begin{quote}
Before the new one: say the Marwadi for a promise, then the Marwadi for trouble.
\end{quote}

\subsection*{You'll want to know: \mw{पुकार}}
\textbf{\mw{पुकार}} — \emph{pukār} — \textquotedblleft{}a cry for help\textquotedblright{}.

\begin{quote}
\mw{हूं} \mw{पुकार} \mw{सुणूं} \mw{हूं।}
\end{quote}

\begin{quote}
\emph{hū̃ pukār suṇū̃ hū̃.}
\end{quote}

\begin{quote}
I hear a cry for help.
\end{quote}

A feminine noun.

\begin{cousinweb}
\textbf{\mw{पुकार}} (\emph{pukār}) is inherited from an imitative \emph{*pukkār}, \textquotedblleft{}shout\textquotedblright{}; CDIAL lists Old Marwari \emph{pukāra}, \textquotedblleft{}complaint\textquotedblright{} (CDIAL 8246).
\end{cousinweb}

\begin{grammarlens}[title={What you've built}]
One more word for asking and answering.
\end{grammarlens}

\subsection*{Your turn}
\begin{itemize}
\raggedright
  \item \emph{Say it:} \emph{pukār}
  \item \emph{Say it:} \emph{pukār}, once more
  \item \emph{Say it:} \emph{hū̃ pukār suṇū̃ hū̃.} — I hear a cry for help.
  \item \emph{Recall:} say the Marwadi for a promise, then the Marwadi for trouble, then say \emph{pukār} again
\end{itemize}

\begin{tcolorbox}[breakable,colback=teal!4,colframe=teal!35!black,title={Before you move on}]
What does \mw{पुकार} mean? (\textquotedblleft{}A cry for help\textquotedblright{}.) Say it once more.
\end{tcolorbox}
\section[\texorpdfstring{\mw{सहारो} (sahāro)}{सहारो (sahāro)}]{\mw{सहारो} (sahāro) — support}

\label{lesson:MW-C274-saharo}



\begin{quote}
Before the new one: say the Marwadi for trouble, then the Marwadi for a cry for help.
\end{quote}

\subsection*{You'll want to know: \mw{सहारो}}
\textbf{\mw{सहारो}} — \emph{sahāro} — \textquotedblleft{}support\textquotedblright{}.

\begin{quote}
\mw{मां} \mw{म्हारो} \mw{सहारो} \mw{है।}
\end{quote}

\begin{quote}
\emph{mā̃ mhāro sahāro hai.}
\end{quote}

\begin{quote}
Mother is my support.
\end{quote}

A masculine noun.

\begin{cousinweb}
\textbf{\mw{सहारो}} (\emph{sahāro}) is inherited from Sanskrit \emph{sahakāra}, \textquotedblleft{}assistance\textquotedblright{}; CDIAL gives no Marwari form in this entry, but lists Hindi \emph{sahārā}, \textquotedblleft{}help\textquotedblright{} (CDIAL 13299).
\end{cousinweb}

\begin{grammarlens}[title={What you've built}]
One more word for asking and answering.
\end{grammarlens}

\subsection*{Your turn}
\begin{itemize}
\raggedright
  \item \emph{Say it:} \emph{sahāro}
  \item \emph{Say it:} \emph{sahāro}, once more
  \item \emph{Say it:} \emph{mā̃ mhāro sahāro hai.} — Mother is my support.
  \item \emph{Recall:} say the Marwadi for trouble, then the Marwadi for a cry for help, then say \emph{sahāro} again
\end{itemize}

\begin{tcolorbox}[breakable,colback=teal!4,colframe=teal!35!black,title={Before you move on}]
What does \mw{सहारो} mean? (\textquotedblleft{}Support\textquotedblright{}.) Say it once more.
\end{tcolorbox}
\section[\texorpdfstring{\mw{दया} (dayā)}{दया (dayā)}]{\mw{दया} (dayā) — mercy}

\label{lesson:MW-C274-daya}



\begin{quote}
Before the new one: say the Marwadi for a cry for help, then the Marwadi for support.
\end{quote}

\subsection*{You'll want to know: \mw{दया}}
\textbf{\mw{दया}} — \emph{dayā} — \textquotedblleft{}mercy\textquotedblright{}.

\begin{quote}
\mw{दया} \mw{करो} \mw{सा।}
\end{quote}

\begin{quote}
\emph{dayā karo sā.}
\end{quote}

\begin{quote}
Have mercy, please.
\end{quote}

A feminine noun.

\begin{cousinweb}
\textbf{\mw{दया}} (\emph{dayā}) is a tatsama, taken from Sanskrit \emph{dayā}, \textquotedblleft{}sympathy, compassion\textquotedblright{} (Monier-Williams). \textbf{\mw{दयाळु}}, kind, is built on it.
\end{cousinweb}

\begin{grammarlens}[title={What you've built}]
One more word for asking and answering.
\end{grammarlens}

\subsection*{Your turn}
\begin{itemize}
\raggedright
  \item \emph{Say it:} \emph{dayā}
  \item \emph{Say it:} \emph{dayā}, once more
  \item \emph{Say it:} \emph{dayā karo sā.} — Have mercy, please.
  \item \emph{Recall:} say the Marwadi for a cry for help, then the Marwadi for support, then say \emph{dayā} again
\end{itemize}

\begin{tcolorbox}[breakable,colback=teal!4,colframe=teal!35!black,title={Before you move on}]
What does \mw{दया} mean? (\textquotedblleft{}Mercy\textquotedblright{}.) Say it once more.
\end{tcolorbox}
\section[\texorpdfstring{\mw{भीख} (bhīkh)}{भीख (bhīkh)}]{\mw{भीख} (bhīkh) — alms}

\label{lesson:MW-C274-bhikh}



\begin{quote}
Before the new one: say the Marwadi for support, then the Marwadi for mercy.
\end{quote}

\subsection*{You'll want to know: \mw{भीख}}
\textbf{\mw{भीख}} — \emph{bhīkh} — \textquotedblleft{}alms\textquotedblright{}.

\begin{quote}
\mw{वो} \mw{भीख} \mw{मांगै} \mw{है।}
\end{quote}

\begin{quote}
\emph{vo bhīkh māṅgai hai.}
\end{quote}

\begin{quote}
He begs for alms.
\end{quote}

A feminine noun.

\begin{cousinweb}
\textbf{\mw{भीख}} (\emph{bhīkh}) is inherited from Sanskrit \emph{bhikṣā}, \textquotedblleft{}alms\textquotedblright{}; CDIAL lists Old Marwari \emph{bhīkha}, \textquotedblleft{}alms\textquotedblright{} (CDIAL 9485).
\end{cousinweb}

\begin{grammarlens}[title={What you've built}]
One more word for asking and answering.
\end{grammarlens}

\subsection*{Your turn}
\begin{itemize}
\raggedright
  \item \emph{Say it:} \emph{bhīkh}
  \item \emph{Say it:} \emph{bhīkh}, once more
  \item \emph{Say it:} \emph{vo bhīkh māṅgai hai.} — He begs for alms.
  \item \emph{Recall:} say the Marwadi for support, then the Marwadi for mercy, then say \emph{bhīkh} again
\end{itemize}

\begin{tcolorbox}[breakable,colback=teal!4,colframe=teal!35!black,title={Before you move on}]
What does \mw{भीख} mean? (\textquotedblleft{}Alms\textquotedblright{}.) Say it once more.
\end{tcolorbox}
